% generated by requirements_tex.py from pacman-firmware ghdl/agents/requirements.txt:  do not edit
{\small
\begin{longtable}{@{}>{\raggedright\arraybackslash}p{0.09\linewidth}>{\raggedright\arraybackslash}p{0.07\linewidth}>{\raggedright\arraybackslash}p{0.77\linewidth}@{}}
\toprule
ID & Where & Requirement \\
\midrule
\endhead
\addlinespace[4pt]\multicolumn{3}{@{}l}{\textbf{REG:  Register access (all units)}} \\
\texttt{REG-1} & \multicolumn{2}{>{\raggedright\arraybackslash}p{\dimexpr0.84\linewidth+2\tabcolsep}@{}}{\emph{Every register is reachable at its address, and nothing else answers.}} \\
\texttt{REG-1.1} & sim+hw & Each unit acknowledges reads and writes of its own defined registers. \\
\texttt{REG-1.2} & sim+hw & Undefined addresses, other units' scopes, writes to read-only registers and reads of the broadcast channel are not acknowledged, and change nothing. \\
\texttt{REG-1.3} & sim+hw & A request nobody acknowledges is completed by the bridge after its timeout (a read returns 0xEEEEEEEE), and the bus keeps working. \\
\texttt{REG-2} & \multicolumn{2}{>{\raggedright\arraybackslash}p{\dimexpr0.84\linewidth+2\tabcolsep}@{}}{\emph{Registers hold what is written, independently.}} \\
\texttt{REG-2.1} & sim+hw & Every read/write register reads back every value written to it, all bits. \\
\texttt{REG-2.2} & sim+hw & A write changes only the register addressed: no crosstalk between registers, channels or units. \\
\texttt{REG-2.3} & sim+hw & A broadcast write (channel 0x3B) sets the register of all 40 UARTs. \\
\texttt{REG-3} & \multicolumn{2}{>{\raggedright\arraybackslash}p{\dimexpr0.84\linewidth+2\tabcolsep}@{}}{\emph{Register access timing.}} \\
\texttt{REG-3.1} & sim & Global and Timing answer reads and writes one clock after the request; RX and TX answer writes one clock and reads two clocks after the request. \\
\texttt{REG-3.2} & sim & Back-to-back requests, to one unit or across units and clock domains, are answered in order, each once, with no mixed data. \\
\texttt{REG-3.3} & sim & A read in the same clock as a write of the same register returns the old value (current behaviour, no stated contract). \\
\texttt{REG-4} & \multicolumn{2}{>{\raggedright\arraybackslash}p{\dimexpr0.84\linewidth+2\tabcolsep}@{}}{\emph{Reset.}} \\
\texttt{REG-4.1} & sim & Reset restores every register to its default, in every unit. \\
\addlinespace[4pt]\multicolumn{3}{@{}l}{\textbf{GLB:  Global unit}} \\
\texttt{GLB-1} & \multicolumn{2}{>{\raggedright\arraybackslash}p{\dimexpr0.84\linewidth+2\tabcolsep}@{}}{\emph{The board controls do exactly what the register says.}} \\
\texttt{GLB-1.1} & sim & ENABLES bit 16 drives TILE\_POWER\_ENABLE and bits 9-0 TILE\_ENABLES; the other bits drive nothing. \\
\texttt{GLB-1.2} & sim & LEDS bits 1-0 drive the two PL LEDs; the other bits drive nothing. \\
\texttt{GLB-1.3} & sim & After reset, tile power, the tiles and the LEDs are off. \\
\texttt{GLB-2} & \multicolumn{2}{>{\raggedright\arraybackslash}p{\dimexpr0.84\linewidth+2\tabcolsep}@{}}{\emph{The firmware identifies itself.}} \\
\texttt{GLB-2.1} & sim+hw & STATUS reads 0x0000CAFE. \\
\texttt{GLB-2.2} & sim+hw & The version registers read the firmware and hardware versions and the build information. \\
\addlinespace[4pt]\multicolumn{3}{@{}l}{\textbf{TX:  TX unit}} \\
\texttt{TX-1} & \multicolumn{2}{>{\raggedright\arraybackslash}p{\dimexpr0.84\linewidth+2\tabcolsep}@{}}{\emph{Every word of a DMA packet is sent on the UARTs of its mask, intact, and nowhere else.}} \\
\texttt{TX-1.1} & sim+hw & The channels of the mask send their own word; the others stay idle; mask bits above 39 are ignored. \\
\texttt{TX-1.2} & sim+hw & All 40 channels can send at once, each its own word. \\
\texttt{TX-1.3} & sim+hw & Back-to-back packets are sent in order, with no word dropped or repeated; the next packet is read while the UARTs send. \\
\texttt{TX-2} & \multicolumn{2}{>{\raggedright\arraybackslash}p{\dimexpr0.84\linewidth+2\tabcolsep}@{}}{\emph{The UART frames are as configured.}} \\
\texttt{TX-2.1} & sim+hw & A frame is a start bit, the 64 data bits least significant first, and a stop bit, one bit per BAUD period. \\
\texttt{TX-2.2} & sim & DELAY adds exactly that many idle bit periods between frames. \\
\texttt{TX-2.3} & sim & PHASE moves the bit edges within the BAUD period. \\
\texttt{TX-3} & \multicolumn{2}{>{\raggedright\arraybackslash}p{\dimexpr0.84\linewidth+2\tabcolsep}@{}}{\emph{The TX unit can be monitored.}} \\
\texttt{TX-3.1} & sim+hw & STARTS and BEATS count each channel's frames and hand-overs (equal in normal operation); ZERO\_CNTS clears them. \\
\texttt{TX-3.2} & sim+hw & LOOK shows the word last loaded for the selected channel. \\
\texttt{TX-3.3} & sim & The status registers show the channel and buffer states. \\
\addlinespace[4pt]\multicolumn{3}{@{}l}{\textbf{RX:  RX unit}} \\
\texttt{RX-1} & \multicolumn{2}{>{\raggedright\arraybackslash}p{\dimexpr0.84\linewidth+2\tabcolsep}@{}}{\emph{Every word received reaches the stream once, intact, with its header and timestamp.}} \\
\texttt{RX-1.1} & sim+hw & A word on one channel comes out as header, timestamp and data, followed by the trailer. \\
\texttt{RX-1.2} & sim+hw & Words on several channels at once come out in channel order, in one packet. \\
\texttt{RX-1.3} & sim+hw & Every data bit arrives as sent. \\
\texttt{RX-1.4} & sim+hw & The header carries the word type (WORD\_TYPE\_LUT, by data bits 1-0), the PACMAN id and the channel id. \\
\texttt{RX-2} & \multicolumn{2}{>{\raggedright\arraybackslash}p{\dimexpr0.84\linewidth+2\tabcolsep}@{}}{\emph{The receiver copes with real lines.}} \\
\texttt{RX-2.1} & sim & A start bit at any phase of the BAUD period is received (MODE 1). \\
\texttt{RX-2.2} & sim & A glitch shorter than a bit is outvoted, and counted as noise. \\
\texttt{RX-2.3} & sim & A bad stop bit is a frame error: counted, no word. \\
\texttt{RX-2.4} & sim+hw & MODE 0 turns a channel off; INPUT selects the line (idle, PISO, loopback, test pattern). \\
\texttt{RX-3} & \multicolumn{2}{>{\raggedright\arraybackslash}p{\dimexpr0.84\linewidth+2\tabcolsep}@{}}{\emph{The stream is cut into DMA packets as configured, and congestion is accounted for.}} \\
\texttt{RX-3.1} & sim & When the stream stalls, waiting words are held; a further word on a channel still holding one is lost, and counted. \\
\texttt{RX-3.2} & sim+hw & BUFFER\_CONFIG's maximum closes a packet after that many words. \\
\texttt{RX-3.3} & sim+hw & BUFFER\_CONFIG's timeout lets later words join the open packet, then closes it. \\
\texttt{RX-4} & \multicolumn{2}{>{\raggedright\arraybackslash}p{\dimexpr0.84\linewidth+2\tabcolsep}@{}}{\emph{The extra channels add their words when enabled.}} \\
\texttt{RX-4.1} & sim+hw & Heartbeat: a word every HEARTBEAT\_CONFIG clocks (HEADER\_A, timestamp, zero data). \\
\texttt{RX-4.2} & sim & Rollover: a word when the timestamp equals ROLLOVER\_CONFIG (HEADER\_B). \\
\texttt{RX-4.3} & sim+hw & Markers 2 and 3: a word for each marker pulse (HEADER\_C, HEADER\_D). \\
\texttt{RX-5} & \multicolumn{2}{>{\raggedright\arraybackslash}p{\dimexpr0.84\linewidth+2\tabcolsep}@{}}{\emph{The test pattern generator feeds any channel.}} \\
\texttt{RX-5.1} & sim+hw & A channel on the test pattern receives the pattern word, frame after frame; the generator's status and counters follow. \\
\texttt{RX-6} & \multicolumn{2}{>{\raggedright\arraybackslash}p{\dimexpr0.84\linewidth+2\tabcolsep}@{}}{\emph{The RX unit can be monitored.}} \\
\texttt{RX-6.1} & sim+hw & The counters (STARTS, BEATS, UPDATES, LOST, FRAME\_ERR, NOISE) count their events; ZERO\_CNTS clears them. \\
\texttt{RX-6.2} & sim+hw & LOOK shows the last word received on the selected channel. \\
\texttt{RX-6.3} & sim & The status registers show the channel and buffer states. \\
\texttt{RX-6.4} & sim & FIFO\_CNT shows the FIFO count, FIFO\_MAX its maximum since ZERO\_CNTS. \\
\addlinespace[4pt]\multicolumn{3}{@{}l}{\textbf{ATC:  Timing unit}} \\
\texttt{ATC-1} & \multicolumn{2}{>{\raggedright\arraybackslash}p{\dimexpr0.84\linewidth+2\tabcolsep}@{}}{\emph{The clocks and the timestamp run as configured.}} \\
\texttt{ATC-1.1} & sim & BAUD and UCLK follow CONFIG\_BAUD and CONFIG\_UART. \\
\texttt{ATC-1.2} & sim+hw & The timestamp counts UART periods, and reaches the RX unit. \\
\texttt{ATC-1.3} & sim+hw & Marker 0 (SYNC) restarts the timestamp. \\
\texttt{ATC-2} & \multicolumn{2}{>{\raggedright\arraybackslash}p{\dimexpr0.84\linewidth+2\tabcolsep}@{}}{\emph{Each stimulus goes where its destination register sends it.}} \\
\texttt{ATC-2.1} & sim & A poke reaches G and/or H on the tiles of its destination and of its own mask, for its duration. \\
\texttt{ATC-2.2} & sim+hw & A stimulus reaches the markers of its destination. \\
\texttt{ATC-2.3} & sim & LEMO A and B, and every poke, route by their own destination registers. \\
\texttt{ATC-2.4} & sim & LOGIC A and B never fire (disabled in this version). \\
\texttt{ATC-3} & \multicolumn{2}{>{\raggedright\arraybackslash}p{\dimexpr0.84\linewidth+2\tabcolsep}@{}}{\emph{The G and H outputs are shaped as configured.}} \\
\texttt{ATC-3.1} & sim & CONFIG\_G and CONFIG\_H set the phase of the outputs within the UART period, and invert tiles. \\
\texttt{ATC-4} & \multicolumn{2}{>{\raggedright\arraybackslash}p{\dimexpr0.84\linewidth+2\tabcolsep}@{}}{\emph{The LEMO inputs give one stimulus per clean edge.}} \\
\texttt{ATC-4.1} & sim & A rising edge held two clocks or more is one stimulus; glitches and bounces are not. \\
\texttt{ATC-5} & \multicolumn{2}{>{\raggedright\arraybackslash}p{\dimexpr0.84\linewidth+2\tabcolsep}@{}}{\emph{The counters count, and are controlled by their commands.}} \\
\texttt{ATC-5.1} & sim+hw & READ latches the edges and the integral of the chosen signal; START, STOP and CLEAR act on all. \\
\bottomrule
\end{longtable}
}
